\begin{proof}[Proof of \cref{obj:thm:sharp-annotation-frontier}]
Fix the public parameters as in the statement. Throughout the proof, use
\[
N=n+m,\qquad S=\alpha+\beta,\qquad
\Delta=2\gamma+d+\frac{\gamma d}{S},\qquad
S_{\mathrm{crit}}=\frac{\gamma d}{2\gamma+d},\qquad
q_*=\frac{\gamma d}{S(2\gamma+d)},
\]
and
\[
r_o(n)=n^{-\gamma/(2\gamma+d)},\qquad
r_*=\max\{r_o(n),(nN)^{-\gamma/\Delta}\},
\]
where the last equality is \cref{obj:def:sharp-rate}. All sample sizes below satisfy \(n,m\in\mathbb N\) and \(n\ge2\).

\begin{enumerate}
\item \emph{Supplied-propensity lower bound.}
By \cref{obj:thm:supplied-propensity}, there is a constant \(c_O>0\), depending on the fixed public parameters, such that
\[
c_O r_o(n)\le R^e(n,m)
\qquad\text{for every }n,m.
\]
Every admissible decision \(T\) induces a supplied-propensity decision by
\[
T^e(\mathcal D_{n,m},U,e)=T(\mathcal D_{n,m},U)
\]
for every function \(e:[0,1]^d\to\mathbb R\). For each supplied function this decision is Borel, and its loss under each \(P\in\mathcal M\) is
\[
\int
\left|T^e(\mathcal D_{n,m},U,e_P)-\tau_P(x_0)\right|
\,d\mathsf E_{n,m}(P)
=
\mathcal R_{n,m}(T,P).
\]
Taking the supremum over the same primitive class and then the infimum over decisions, using \cref{obj:def:oracle-risk,obj:def:minimax}, gives
\[
R^e(n,m)\le R(n,m).
\]
Consequently,
\[
c_O r_o(n)\le R(n,m)
\qquad\text{for every }n,m.
\]
% lean: CausalSmith.Stat.TwosamplePointcateAnnotationFrontier.oracleRisk_le_minimaxRisk

\item \emph{A uniform lower constant.}
First suppose \(S<S_{\mathrm{crit}}\). By
\cref{obj:lem:nuisance-frontier-converse}, there is a constant \(c_B>0\), depending on the fixed public parameters, such that
\[
c_B(nN)^{-\gamma/\Delta}\le R(n,m)
\qquad\text{whenever }n\le N\le n^{q_*}.
\]
Choose
\[
c=\min\{c_O,c_B\}>0.
\]

We establish the rate comparisons used in this case and in the remaining smoothness case. The parameter conditions imply
\[
S>0,\qquad \gamma>0,\qquad 2\gamma+d>0,\qquad
\Delta>0,\qquad q_*>0,\qquad N\ge n\ge2.
\]
Direct substitution gives
\[
q_*=
\frac{\gamma d}{S(2\gamma+d)}
=
\frac{S_{\mathrm{crit}}}{S},
\]
and
\[
\Delta
=
(2\gamma+d)+\frac{\gamma d}{S}
=
(2\gamma+d)
\left(1+\frac{\gamma d}{S(2\gamma+d)}\right)
=
(2\gamma+d)(1+q_*).
\]
% lean: CausalSmith.Stat.TwosamplePointcateAnnotationFrontier.rate_exponent_identities

Since \(n>0\), the laws of real powers therefore give the boundary identity
\[
\begin{aligned}
(n\,n^{q_*})^{-\gamma/\Delta}
&=(n^{1+q_*})^{-\gamma/\Delta}\\
&=n^{-(1+q_*)\gamma/\Delta}\\
&=n^{-\gamma/(2\gamma+d)}
=r_o(n).
\end{aligned}
\]
The exponent \(-\gamma/\Delta\) is negative. Thus, when \(N\ge n^{q_*}\),
\[
0<n\,n^{q_*}\le nN
\quad\Longrightarrow\quad
(nN)^{-\gamma/\Delta}
\le(n\,n^{q_*})^{-\gamma/\Delta}
=r_o(n).
\]
When \(N\le n^{q_*}\), instead,
\[
0<nN\le n\,n^{q_*}
\quad\Longrightarrow\quad
r_o(n)
=(n\,n^{q_*})^{-\gamma/\Delta}
\le(nN)^{-\gamma/\Delta}.
\]
For completeness, if \(S\ge S_{\mathrm{crit}}\), then
\[
q_*=\frac{S_{\mathrm{crit}}}{S}\le1,
\qquad
n^{q_*}\le n^1=n\le N,
\]
where the power comparison uses \(n\ge1\). Taking the maximum in the definition of \(r_*\) consequently gives
\[
\begin{aligned}
S\ge S_{\mathrm{crit}}
&\quad\Longrightarrow\quad r_*=r_o(n),\\
S<S_{\mathrm{crit}},\ N\le n^{q_*}
&\quad\Longrightarrow\quad r_*=(nN)^{-\gamma/\Delta},\\
S<S_{\mathrm{crit}},\ N\ge n^{q_*}
&\quad\Longrightarrow\quad r_*=r_o(n).
\end{aligned}
\]
The boundary identity also gives
\[
N=n^{q_*}
\quad\Longrightarrow\quad
(nN)^{-\gamma/\Delta}=r_o(n).
\]
% lean: CausalSmith.Stat.TwosamplePointcateAnnotationFrontier.rate_branches

Return to the case \(S<S_{\mathrm{crit}}\). If \(N\le n^{q_*}\), then \(N\ge n\), so the nuisance lower bound applies and yields
\[
c r_*
=c(nN)^{-\gamma/\Delta}
\le c_B(nN)^{-\gamma/\Delta}
\le R(n,m).
\]
If \(N>n^{q_*}\), the supplied-propensity lower bound yields
\[
c r_*=c r_o(n)\le c_O r_o(n)\le R(n,m).
\]
Here the comparisons of constants preserve the inequalities because both rate terms are nonnegative.

In the remaining case \(S\ge S_{\mathrm{crit}}\), choose
\[
c=c_O>0.
\]
Then \(r_*=r_o(n)\), and hence
\[
c r_*=c_O r_o(n)\le R(n,m).
\]
This constructs a positive constant \(c\), depending on the fixed public parameters and valid simultaneously for every admissible sample-size pair, such that
\[
c r_*\le R(n,m).
\]
% lean: CausalSmith.Stat.TwosamplePointcateAnnotationFrontier.sharpRate_minimax_lower

\item \emph{Attainment and the upper constant.}
By \cref{obj:thm:rectangular-upper}, there is a real constant \(C_U>0\), depending on the fixed public parameters, such that the publicly tuned decision is Borel measurable and its risk under every \(P\in\mathcal M\) satisfies
\[
\mathcal R_{n,m}(\widehat T_{\mathrm{up}},P)
\le C_U r_{\mathrm{up}}(n,m).
\]
By \cref{obj:def:sharp-decision,obj:def:sharp-rate},
\[
T_*=\widehat T_{\mathrm{up}},
\qquad
r_{\mathrm{up}}(n,m)=r_*.
\]
Set
\[
C=\max\{c,C_U\}.
\]
Then \(0<c\le C<\infty\). For every admissible sample-size pair, \(T_*\) is Borel measurable, and taking the supremum of the pointwise risk bound gives
\[
\sup_{P\in\mathcal M}\mathcal R_{n,m}(T_*,P)
\le C_U r_*
\le C r_*.
\]
The last inequality follows from \(C_U\le C\) and \(r_*\ge0\).
% lean: CausalSmith.Stat.TwosamplePointcateAnnotationFrontier.sharp_annotation_frontier

\item \emph{Minimax comparison and arbitrary decisions.}
Since \(T_*\) belongs to the decision class in \cref{obj:def:minimax},
\[
R(n,m)\le
\sup_{P\in\mathcal M}\mathcal R_{n,m}(T_*,P).
\]
Together with the bounds already established, this proves
\[
c r_*
\le R(n,m)
\le\sup_{P\in\mathcal M}\mathcal R_{n,m}(T_*,P)
\le C r_*.
\]
For any Borel real-valued decision \(T\) using the original dataset and the independent uniform randomizer, the same infimum comparison gives
\[
c r_*\le R(n,m)
\le\sup_{P\in\mathcal M}\mathcal R_{n,m}(T,P).
\]
The rate comparisons and boundary identity established above supply all three stated branches and their equality at \(N=n^{q_*}\).
% lean: CausalSmith.Stat.TwosamplePointcateAnnotationFrontier.sharp_annotation_frontier
\end{enumerate}
\end{proof}
